\documentclass[12pt,a4paper]{article}
\usepackage[margin=25mm, headheight=15pt, headsep=10pt]{geometry}
\usepackage{fontspec}
\usepackage[table]{xcolor} % Note: table option is required for rowcolors
\usepackage{titlesec}
\usepackage{tocloft}
\usepackage{fancyhdr}
\usepackage{booktabs}
\usepackage{tcolorbox}
\tcbuselibrary{skins, breakable}
\usepackage[colorlinks=true, linkcolor=emerald, urlcolor=emerald, bookmarks=true]{hyperref}

\definecolor{emerald}{RGB}{0,102,51}
\definecolor{darkred}{RGB}{139,0,0}
\definecolor{lightgray}{RGB}{245,245,245}

\usepackage[nil]{babel}
\babelprovide[import, main]{bengali}
\babelprovide[import]{arabic}
\babelfont{rm}[Renderer=HarfBuzz,Scale=1.0]{Noto Serif Bengali}
\babelfont{sf}[Renderer=HarfBuzz]{Noto Sans Bengali}
\babelfont{tt}[Renderer=HarfBuzz]{Noto Sans Bengali}
\babelfont[arabic]{rm}[Renderer=HarfBuzz,Scale=1.1]{Noto Naskh Arabic}

% Section styling
\titleformat{\section}{\Large\bfseries\color{emerald}}{\thesection.}{0.5em}{}
\titleformat{\subsection}{\large\bfseries\color{emerald!85!black}}{\thesubsection.}{0.5em}{}
\titleformat{\subsubsection}{\normalsize\bfseries\color{emerald!70!black}}{\thesubsubsection.}{0.5em}{}
\titlespacing*{\section}{0pt}{18pt}{8pt}

% Fancy Header/Footer setup
\pagestyle{fancy}
\fancyhf{} % clear all header and footer fields
\fancyhead[L]{\color{emerald}\small\textbf{\nouppercase{\leftmark}}}
\fancyfoot[L]{\scriptsize\color{black!50}{\lat{AI}}-সংকলিত, আলেম-যাচাইকৃত নয়}
\fancyfoot[C]{\color{emerald!70!black}\thepage}
\renewcommand{\headrulewidth}{0.4pt}
\renewcommand{\headrule}{\hbox to\headwidth{\color{emerald}\leaders\hrule height \headrulewidth\hfill}}
\renewcommand{\footrulewidth}{0pt}

% Custom boxes
% 1. Ayat/Hadith Box
\newtcolorbox{ayatbox}[1][]{
    enhanced, breakable, 
    colback=emerald!5, 
    colframe=emerald, 
    boxrule=0.5pt, 
    arc=3pt, 
    left=10pt, right=10pt, top=10pt, bottom=10pt,
    #1
}

% 2. Quote/Translation Box
\newtcolorbox{quotebox}[1][]{
    enhanced, breakable,
    frame hidden,
    colback=lightgray,
    borderline west={3pt}{0pt}{emerald},
    left=15pt, right=10pt, top=10pt, bottom=10pt,
    sharp corners,
    #1
}

% 3. AI-generated notice box
\newtcolorbox{warnbox}[1][]{
    enhanced, breakable,
    colback=red!4,
    colframe=darkred,
    boxrule=0.8pt,
    arc=3pt,
    left=10pt, right=10pt, top=10pt, bottom=10pt,
    #1
}

% Commands
\newcommand{\arab}[1]{\foreignlanguage{arabic}{#1}}
\newcommand{\kib}[1]{\textcolor{emerald}{\textbf{#1}}}

% Latin (English) fallback: Noto Serif Bengali has NO Latin glyphs
\newfontfamily\latfont[Renderer=HarfBuzz]{Noto Serif}
\newcommand{\lat}[1]{{\latfont #1}}

\setlength{\parskip}{5pt}
\setlength{\parindent}{0pt}

% Fix for missing bullet in Noto Serif Bengali
\renewcommand{\labelitemi}{$\bullet$}



\begin{document}
\begin{center}{\Large\bfseries\color{emerald} \lat{09}-জালসা-দোয়া}\end{center}
\vspace{1em}
\section*{০৯ — জালসা (দুই সিজদার মাঝে বৈঠক): \textit{\lat{Rabbighfir} \lat{li}}}

\begin{warnbox}
 \textbf{গুরুত্বপূর্ণ নোটিশ:} এই কন্টেন্টটি \lat{AI} দিয়ে প্রস্তুত। কেবল সহীহ/হাসান হাদিস রাখা হয়েছে। আলেম-যাচাইকৃত নয়।
\end{warnbox}

\medskip\hrule\medskip

\subsection*{১. সূত্র কার্ড}

\begin{table}[h]\centering\small
\begin{tabular}{ll}
বিষয় & বিবরণ \\
\hline
\textbf{বিষয়} & দুই সিজদার মাঝে বসে বলা দোয়া \\
\textbf{সূত্র} & সহীহ মুসলিম ৫৭৯-৫৮০, সুনানে আবু দাউদ ৮৫০ \\
\textbf{মর্যাদা} & \textbf{সুন্নাতে মুআক্কাদাহ} \\
\hline
\end{tabular}
\end{table}

\medskip\hrule\medskip

\subsection*{২. মূল আরবি}

\subsubsection*{\textbf{প্রথম দোয়া:}}
\begin{center}\arab{رَبِّ اغْفِرْ لِي رَبِّ اغْفِرْ لِي}\end{center}

\subsubsection*{\textbf{দীর্ঘ দোয়া (বিকল্প):}}
\begin{center}\arab{رَبِّ اغْفِرْ لِي وَارْحَمْنِي وَاهْدِنِي وَعَافِنِي وَارْزُقْنِي}\end{center}

\medskip\hrule\medskip

\subsection*{৩. বাংলা উচ্চারণ}

\textbf{প্রথম:} রাব্বিগফির লী, রাব্বিগফির লী

\textbf{দীর্ঘ:} রাব্বিগফির লী, ওয়ারহামনী, ওয়াহদিনী, ওয়া'আফিনী, ওয়ারযুকনী

\medskip\hrule\medskip

\subsection*{৪. শব্দ-শব্দ অর্থ}

\begin{table}[h]\centering\small
\begin{tabular}{lll}
আরবি & বাংলা & \lat{English} \\
\hline
\textbf{\arab{رَبِّ}} & হে আমার রব & \textit{\lat{O} \lat{my} \lat{Lord}} \\
\textbf{\arab{اغْفِرْ} \arab{لِي}} & আমাকে ক্ষমা কর & \textit{\lat{forgive} \lat{me}} \\
\textbf{\arab{وَارْحَمْنِي}} & ও আমার উপর দয়া কর & \textit{\lat{and} \lat{have} \lat{mercy} \lat{on} \lat{me}} \\
\textbf{\arab{وَاهْدِنِي}} & ও আমাকে হেদায়াত দাও & \textit{\lat{and} \lat{guide} \lat{me}} \\
\textbf{\arab{وَعَافِنِي}} & ও আমাকে আরোগ্য দাও & \textit{\lat{and} \lat{grant} \lat{me} \lat{well-being}} \\
\textbf{\arab{وَارْزُقْنِ}ি} & ও আমাকে রিযিক দাও & \textit{\lat{and} \lat{provide} \lat{for} \lat{me}} \\
\hline
\end{tabular}
\end{table}

\medskip\hrule\medskip

\subsection*{৫. পূর্ণ বাংলা অনুবাদ}

\textbf{\lat{“}হে আমার রব! আমাকে ক্ষমা কর।\lat{”}}

\textbf{\lat{“}হে আমার রব! আমাকে ক্ষমা কর, আমার উপর দয়া কর, আমাকে হেদায়াত দাও, আমাকে আরোগ্য দাও ও রিযিক দাও।\lat{”}}

\medskip\hrule\medskip

\subsection*{৬. সংক্ষিপ্ত ব্যাখ্যা}

\subsubsection*{\textbf{জালসা কী?}}
\textbf{\lat{Jalsah} (জালসা)} মানে প্রথম সিজদা থেকে উঠে বসে ওয়াকফা করা — যতক্ষণ "রাব্বিগফির লী" বলা যায়। এরপর দ্বিতীয় সিজদায় যান।

\subsubsection*{\textbf{"রাব্বিগফির লী" — ক্ষমা প্রার্থনা}}
সিজদায় "আ'লা" বলার পর জালসায় "ক্ষমা" চাওয়া — এটি \textbf{\lat{repentance} (তাওবা)} ও \textbf{\lat{humility} (বিনয়)}। বান্দা আল্লাহর সামনে নিজেকে \textbf{\lat{small} (ছোট)} ভাবে দেখছে।

\subsubsection*{\textbf{দীর্ঘ দোয়া — ৫টি প্রার্থনা}}
১। \textbf{\lat{Ighfir} \lat{li} (ক্ষমা)} — গুনাহ মাফ
২। \textbf{\lat{Rahamni} (রহমাত)} — দয়া
৩। \textbf{\lat{Hudini} (হেদায়াত)} — পথ দেখানো
৪। \textbf{'\lat{Aafini} (আরোগ্য)} — শারীরিক-মানসিক সুস্থতা
৫। \textbf{\lat{Rzuqni} (রিযিক)} — জীবিকা

\medskip\hrule\medskip

\subsection*{৭. খুশুর জন্য মূল পয়েন্টস}

\begin{table}[h]\centering\small
\begin{tabular}{ll}
পয়েন্ট & কীভাবে প্রয়োগ করবেন \\
\hline
\textbf{ধীরতা} & দোয়াটি ধীরে, স্পষ্ট করে বলুন \\
\textbf{অর্থ} & প্রতিটি শব্দের অর্থ মনে করুন \\
\textbf{বিশ্বাস} & আল্লাহ শুনছেন ও কবুল করবেন — এই \lat{conviction} (দৃঢ় বিশ্বাস) রাখুন \\
\textbf{পুনরাবৃত্তি} & প্রতিটি সিজদার পর এই দোয়া — এটে \lat{repetition} (পুনরাবৃত্তি) সবচেয়ে ভালো \\
\hline
\end{tabular}
\end{table}

\medskip\hrule\medskip

\subsection*{৮. সংশ্লিষ্ট হাদিস}

\begin{table}[h]\centering\small
\begin{tabular}{ll}
হাদিস & গ্রন্থ \\
\hline
\textbf{\lat{“}সিজদার মাঝে 'রাব্বিগফির লী, রাব্বিগফির লী' বলো।\lat{”}} & সহীহ মুসলিম ৫৭৯ \\
\textbf{\lat{“}সিজদার মাঝে দোয়া করো — বাধা নেই।\lat{”}} & সহীহ মুসলিম ৫৮০ \\
\hline
\end{tabular}
\end{table}

\medskip\hrule\medskip

\textit{শেষ আপডেট: আগস্ট ২০২৬}

\end{document}
